% Fresh strict module checks passed; canonical integration gates remain pending.
\section{A matrix product approximation to one soft-truncated interval}
\label{sec:soft_interval_mpo}

For $t>0$ and an integer $m\geq0$, write
$h_{m,t}(x)=t\sum_{j=1}^m(1-e^{-x/t})^j/j$.
This section constructs an approximation to one regional replacement,
including its exterior and scalar shift, for independent parameters $m\geq1$ and $t\geq3m$.
It uses the corrected label estimate, and does not assert a running-time
bound or the unrestricted soft-truncation lemma of
\cite[lines 1262--1268]{Arad2016RigorousRG}.


\begin{lemma}[Weighted sums of operator chains]
    \label{lem:trace_chain_linear_combination}
    \lean{MPOTensor.traceChainOperator_sum_smul_closing,
        MPOTensor.exists_traceChainOperator_common_linearCombination_varying,
        MPOTensor.exists_traceChainOperator_common_linearCombination,
        MPOTensor.exists_traceChainOperator_linearCombination_varying,
        MPOTensor.exists_traceChainOperator_linearCombination,
        MPOTensor.exists_traceChainOperator_add}
    \leanok

    \uses{def:trace_chain_operator}
    For fixed site tensors $A$, the trace contraction
    $\mathcal T(A,\sum_a c_aC_a)=\sum_a c_a\mathcal T(A,C_a)$
    is linear in its closing matrix. Fix site-dependent operator tensors $A_a$ of virtual dimension $D_a$
    and arbitrary fixed closing matrices $Q_a$. There are site tensors $B$
    of dimension $\sum_aD_a$ and fixed matrices $C_a$, before all complex
    weights $c_a$, such that
    \begin{align}
        \sum_a c_a\,\mathcal T(A_a,Q_a)
        =\mathcal T\!\left(B,\sum_a c_aC_a\right).
        \label{eq:trace_chain_common_closing}
    \end{align}
    The physical chain may be empty. On a nonempty physical chain, a finite weighted sum of
    site-dependent matrix product operators with virtual dimensions
    $D_a$ has actual operator tensors of dimension $\sum_aD_a$.
    The coefficients may be arbitrary complex numbers, and zero
    virtual dimensions are permitted. In particular, $s$ summands
    of common dimension $D$ give dimension $sD$, and adding two
    chains of dimensions $D,E$ gives dimension $D+E$.
\end{lemma}
\begin{proof}\leanok
    Put the site tensors on diagonal blocks. Their ordered product
    has the corresponding products on its diagonal blocks. Put the
    coefficients in the same blocks of the closing matrix. Taking
    the trace gives the required weighted sum. Reindex the virtual
    coordinates by a finite interval and absorb the closing matrix
    into the first site tensor when the chain is nonempty. Before this
    last absorption, diagonal caps with only one nonzero block give the
    fixed coefficient matrices $C_a$ in
    (\ref{eq:trace_chain_common_closing}).
\end{proof}

\begin{lemma}[Exact local sum and scalar]
    \label{lem:local_sum_scalar_mpo}
    \lean{TNLean.MPS.AreaLaw.exists_local_sum_add_scalar_mpo}
    \leanok

    \uses{thm:operator_chain_from_all_cut_ranks}
    Let $q\geq1$ and let $h_i$, $0\leq i<N$, act on the designated
    pairs $\{i,i+1\}$ of $N+1$ sites. For every scalar $c$, the
    operator $\sum_i h_i+cI$ has actual site-dependent operator
    tensors of virtual dimension $1+Nq^2$. Hermiticity and
    positivity are unnecessary for this statement.
\end{lemma}
\begin{proof}\leanok
    \uses{thm:operator_chain_from_all_cut_ranks}
    At every cut, each designated pair contributes rank at most
    $q^2$, and the scalar identity contributes at most one.
    Subadditivity gives the bound at all cuts. Apply the constructed
    operator-chain representation.
\end{proof}

\begin{lemma}[Constructed heat approximants after the energy shift]
    \label{thm:shifted_regional_heat_mpo}
    \lean{TNLean.MPS.AreaLaw.exists_shifted_heat_mpo}
    \leanok

    \uses{thm:physical_cluster_label_mpo}
    Let $A=\sum_i h_i$ be an actual nearest-neighbour sum of
    Hermitian terms of norm at most one, and let $\epsilon I\leq A$.
    If $\beta\geq0$, $\ell\geq1$, and $b_\beta=2(e^\beta-1)\leq1$,
    there are actual tensors of dimension $\ell^2q^\ell$ for an
    operator $Q$ satisfying
    \begin{align}
        \|e^{-\beta(A-\epsilon I)}-Q\|
           \leq e^{N^2b_\beta^\ell}-1.
        \label{eq:soft_interval_shifted_heat}
    \end{align}
\end{lemma}
\begin{proof}\leanok
    \uses{thm:physical_cluster_label_mpo}
    Apply the normalized cluster theorem at the lower energy
    $\epsilon$ and multiply its actual tensors by
    $e^{\beta\epsilon}$ through the closing matrix. The exact
    identity $e^{-\beta(A-\epsilon I)}=e^{\beta\epsilon}e^{-\beta A}$
    makes this scalar cancel the normalization in the error bound.
    The scalar shift is kept separate from the local interactions.
\end{proof}

\begin{lemma}[Actual soft tensors at a specified cluster order]
    \label{thm:soft_mpo_cluster_order}
    \lean{TNLean.MPS.AreaLaw.exists_softTruncation_mpo_of_cluster_order}
    \leanok

    \uses{thm:shifted_regional_heat_mpo, lem:trace_chain_linear_combination}
    Let $U=A-\epsilon I\geq0$ have the preceding actual local
    decomposition, and let $m\geq0$, $t>0$, $\ell\geq1$, and
    $b=2(e^{m/t}-1)\leq1$. There are actual operator tensors of
    dimension $(m+1)\ell^2q^\ell$ for an operator $B$ such that
    \begin{align}
        \|h_{m,t}(U)-B\|
          \leq t2^{m+1}(e^{N^2b^\ell}-1).
        \label{eq:soft_interval_soft_cluster_error}
    \end{align}
\end{lemma}
\begin{proof}\leanok
    \uses{thm:shifted_regional_heat_mpo, lem:trace_chain_linear_combination}
    Expand $h_{m,t}$ as the finite heat sum with coefficients
    $c_a$. Construct each heat approximant directly at time $a/t$,
    and combine its actual tensors by diagonal blocks. Every heat
    time is at most $m/t$, so the triangle inequality and the
    absolute coefficient budget $\sum_a|c_a|\leq t2^{m+1}$ give the
    bound. No approximate heat operator is raised to the power $m$.
\end{proof}

\begin{lemma}[A complete soft Hamiltonian at prescribed norm accuracy]
    \label{thm:soft_hamiltonian_mpo_accuracy}
    \lean{TNLean.MPS.AreaLaw.softHeat_base_bounds,
        TNLean.MPS.AreaLaw.exists_softTruncation_mpo_of_accuracy,
        TNLean.MPS.AreaLaw.exists_softTruncatedHamiltonian_mpo_of_accuracy}
    \leanok

    \uses{thm:soft_mpo_cluster_order, lem:local_sum_scalar_mpo,
        lem:trace_chain_linear_combination}
    Let $h_i$, $0\leq i<N$, be Hermitian matrices on a chain of
    $N+1$ sites of local dimension $q\geq1$. Suppose $h_i$ is supported
    on $\{i,i+1\}$ and $\|h_i\|\leq1$, and write $H=\sum_i h_i$.
    Let $m\geq1$, $t\geq3m$, $\xi>0$, and let $I$ be a set of the
    bond labels. Write $A=\sum_{i\in I}h_i$, suppose
    $\epsilon I\leq A$, and put
    \begin{align}
        W&=t2^{m+1},&b&=2(e^{1/3}-1),\\
        \ell&=\max\left\{1,
          \left\lceil\frac{\log\max\{1,N^2\}
                    -\log\log(1+\xi/W)}{-\log b}\right\rceil\right\}.
        \label{eq:soft_interval_cluster_order}
    \end{align}
    Then $0<b<1$. The actual regional soft term $h_{m,t}(A-\epsilon I)$
    has a norm approximation of error at most $\xi$ with virtual
    dimension $(m+1)\ell^2q^\ell$. The full replacement
    \begin{align}
        \widetilde H=H-A+\epsilon I+h_{m,t}(A-\epsilon I)
        \label{eq:soft_interval_full_operator}
    \end{align}
    has actual approximating tensors of dimension
    \begin{align}
        D&=1+Nq^2+(m+1)\ell^2q^\ell,
        \label{eq:soft_interval_bond}\\
        \|\widetilde H-T\|&\leq\xi.
        \label{eq:soft_interval_full_error}
    \end{align}
\end{lemma}
\begin{proof}\leanok
    \uses{thm:soft_mpo_cluster_order, lem:local_sum_scalar_mpo,
        lem:trace_chain_linear_combination}
    The hypotheses give $m/t\leq1/3$. The strict inequality
    $1/3<\log(3/2)$ gives $b<1$, and every required heat time
    lies in $[0,1/3]$.
    The logarithmic ceiling ensures
    $W(e^{N^2b^\ell}-1)\leq\xi$. Apply the preceding constructed
    approximation to the original terms masked by $I$.
    Represent the complementary local sum and $\epsilon I$ exactly
    by the local-sum lemma, and add its tensors. Thus the
    approximation concerns the full replacement, with no omitted
    exterior or scalar term.
\end{proof}

For $0<\xi\leq W$, put $c=-\log b$ and
$E=2\max\{1,N\}^2W/\xi$. The elementary bound
$\log(1+\xi/W)\geq\xi/(2W)$ gives
\begin{align}
    \ell&\leq1+\frac{\log E}{c},\\
    (m+1)\ell^2q^\ell
      &\leq q(m+1)\left(1+\frac{\log E}{c}\right)^2
          E^{\log q/c}.
\end{align}
Here $\log_+u=\max\{0,\log u\}$. Thus, for fixed $q$, the logarithm of the bond is
$O_q(m+\log_+t+\log(N+1)+\log_+(\xi^{-1}))$.
In particular, $m=O(\log(N+2))$, polynomially bounded $t$, and
inverse-polynomial accuracy give a polynomial bond. Polylogarithmic
bounds for $m$, $\log_+t$ and $\log_+(\xi^{-1})$ give a
quasipolynomial bond. These consequences retain the parameter growth;
they do not assert a polynomial bound for arbitrary $m$.
The source's choice $m=4$ is covered whenever $t\geq12$.
This is an existence result with explicit dimensions; a computational
running-time claim and the multi-region spectral argument are not included.
